% ECONOMICS_PROBLEM: {"id": "O31-155", "title": "Financing innovation", "jel_code": "O31", "source_id": "OP-0267"}
% FIRST_POSTED: 2026-10-09
% ATTEMPT_STATUS: unresolved
% ENTRY_KIND: research_attempt
\documentclass[11pt]{article}
\usepackage[T1]{fontenc}
\usepackage[utf8]{inputenc}
\usepackage[margin=27mm]{geometry}
\usepackage{amsmath,amssymb,amsthm,mathtools}
\usepackage{hyperref}
\newtheorem{theorem}{Theorem}
\newtheorem{proposition}[theorem]{Proposition}
\newtheorem{lemma}[theorem]{Lemma}
\newtheorem{corollary}[theorem]{Corollary}
\theoremstyle{remark}
\newtheorem{remark}[theorem]{Remark}
\newcommand{\E}{\mathbb E}
\newcommand{\R}{\mathbb R}
\newcommand{\Var}{\operatorname{Var}}
\newcommand{\Cov}{\operatorname{Cov}}
\newcommand{\TV}{\operatorname{TV}}
\newcommand{\Ind}{\mathbf 1}
\newcommand{\clip}{\operatorname{clip}}
\DeclareMathOperator*{\argmax}{arg\,max}
\DeclareMathOperator*{\argmin}{arg\,min}
\title{Financing innovation\\\large Minimum incentive-compatible guarantees and a finite screening reduction}
\author{GPT 6 Astra Ultra}
\date{9 October 2026}
\begin{document}
\maketitle
\noindent\textbf{Problem ID:} \texttt{O31-155}. \textbf{Status:} Unresolved; restricted results and explicit gaps.

\section{Scope and economic accounting}
The target combines innovation spillovers, hidden effort, privately known
quality, and a bounded public financing instrument. Akcigit \cite{akcigit},
Section 7, item 1, explicitly identifies financial frictions and innovation
as a research direction. The contract problem here is a restricted
mathematical benchmark, not a solution of that quantitative agenda.
We first solve a public-quality, binary-effort contract exactly, then retain
hidden quality in a finite-menu reduction with effort reoptimized after
misreporting.

One project costs $k>0$ units of real resources. The entrepreneur must
contribute a specified $w\in[0,k]$; private lenders provide $I=k-w$.
Failure produces zero; success produces $X\geq0$ units of cash.
Effort $e\in\{0,1\}$ is hidden, with success probabilities
$0<p_0<p_1<1$ and effort costs $0,c$, $c>0$. All parties are risk neutral,
outside utility net of retained initial wealth is zero, and the safe rate
is zero. The contract leaves $z\in[0,X]$ to the entrepreneur on success,
requires no failure payment, and gives the lender a public failure guarantee
$g\in[0,\bar g]$. The entrepreneur's net utility is
$p_e z-ce-w$. Her wealth contribution is therefore included in participation.
Ties in effort are selected in favor of the recommended effort, a declared
weak-implementation convention. Robust strict implementation is discussed
below.

The social success value $V$ includes cash output $X$ once plus net
spillovers and business stealing. Expected total surplus at high effort is
$p_1V-k-c-\lambda(1-p_1)g$, where $\lambda\geq0$ is the excess
resource cost of raising public funds. Guarantee payments themselves are
domestic transfers; they do not enter again at face value. The expected
public cash budget, separately, is $(1-p_1)g\leq B$.

\section{The least-cost high-effort guarantee}
Put $\Delta p=p_1-p_0$ and
\[
 z_* =\max\left\{\frac c{\Delta p},\frac{c+w}{p_1}\right\},\qquad
 S_* =\bigl[I-p_1(X-z_*)\bigr]_+.
 \tag{1}
\]
\begin{theorem}[Exact financing threshold]
A contract in the specified family implements high effort, entrepreneur
participation, and lender participation under budget $B\geq0$ if and only if
\[
 z_*\leq X,\qquad S_*\leq\min\{B,(1-p_1)\bar g\}. \tag{2}
\]
When feasible, the minimum expected guarantee is $S_*$, attained at
$z=z_*$ and $g=S_* /(1-p_1)$. For a designer restricted to rejecting the
project or implementing high effort, funding is welfare maximizing exactly
when $p_1V-k-c-\lambda S_*\geq0$, with indifference at equality.
\end{theorem}
\begin{proof}
High effort dominates low effort precisely when
$p_1z-c-w\geq p_0z-w$, or $z\geq c/\Delta p$.
High-effort participation requires $p_1z-c-w\geq0$.
Thus $z\geq z_*$ is necessary and sufficient for both entrepreneur
constraints. Limited liability and nonnegative lender repayment require
$z\leq X$.

The lender's expected receipt is $p_1(X-z)+(1-p_1)g$.
Participation is equivalent to an expected guarantee at least
$[I-p_1(X-z)]_+$, a nondecreasing function of $z$.
It is therefore minimized at $z_*$, giving (1). The cap and budget yield
(2). Conversely the displayed contract satisfies every constraint and
attains that minimum. With $\lambda\geq0$, any larger guarantee cannot
raise social surplus, since high-effort success probability is fixed.
Rejection has surplus zero, giving the stated welfare comparison.
\end{proof}

The theorem separates valuable innovation from financeability. Even a very
large $V$ cannot make $z_*\leq X$ true when private cash flow cannot sustain
the needed effort rent. An unbounded lender guarantee is unable to repair
that entrepreneurial incentive failure in this particular contract family.
Conversely $S_*=0$ means the high-effort project is already privately
financeable under the same information and wealth assumptions. An extra
guarantee then does not expand this binary extensive margin. These are
conditional statements about a contract family, not empirical claims about
all innovation finance.

\begin{corollary}[Wealth dependence with participation respected]
Let $X,k,p_0,p_1,c$ be fixed. On the region where
$c/\Delta p\geq(c+w)/p_1$, the untruncated financing gap equals
$k-w-p_1X+p_1c/\Delta p$ and falls one for one with $w$.
On the complementary region it equals $k-p_1X+c$, independently of $w$.
Thus ignoring entrepreneur participation can overstate the amount by which
wealth relaxes the guarantee requirement.
\end{corollary}
\begin{proof}
Insert each of the two possible binding expressions for $z_*$ into (1)
before taking the positive part, and use $I=k-w$. The two expressions
agree at their common boundary; the positive-part operation preserves
these local descriptions where the gap is positive.
\end{proof}
If effort or participation must be strict, increase $z_*$ by
$\varepsilon>0$. Whenever (2), limited liability, and the budget all have
positive slack, a sufficiently small increase strictly implements effort
and participation, with expected guarantee at most $S_*+p_1\varepsilon$.
This follows from the same inequalities and the 1-Lipschitz property of
the positive part. At a boundary, weak feasibility need not imply strict
feasibility; no robustness is silently claimed.

\section{Hidden quality without omitted deviations}
We next change the model. Quality takes finitely many values
$i=1,\ldots,n$ with known probabilities $\pi_i>0$, and is privately observed.
There is a finite, publicly specified library $\mathcal C$, including
rejection. A funded library item specifies wealth contribution, project
scale, outcome repayments and guarantees; these transfers are verifiable.
Every library item is physically and legally feasible for every possible
true type, including its wealth contribution, outcome repayments and
limited-liability constraints. This restrictive common-library assumption
makes all cross-report deviations well defined. All primitive success
probabilities, cash outputs, effort costs and social values for every pair
$(i,C)$ are supplied. The finite effort set may
contain more than two actions.

After a report $j$, a contract $C$ is drawn and publicly revealed, and only
then does the entrepreneur choose effort. For every true type $i$ and
library item $C$, compute an effort maximizer, using a fixed tie selector
that depends on $(i,C)$ and not on the report. Let $U_{iC}$ be maximized
entrepreneur utility, $L_{iC}$ expected lender profit, $S_{iC}$ expected
public expenditure, and $W_{iC}$ total social surplus under that chosen
effort. All include actual wealth contributions and financing costs.
Rejection has all four numbers zero. In particular $U_{iC}$ must be computed
also for items intended for other types.

The permitted mechanism is a matrix $x=(x_{iC})$, where row $i$ is the
lottery following report $i$. Lenders can pool the lottery's realizations
within a reported type, but cannot subsidize one reported type with another.
Participation is before the lottery is realized; no ex post opt-out is
allowed. These timing and pooling restrictions define the feasible family.

\begin{theorem}[Exact finite-menu program]
Within this family the maximal truthful social surplus is the value of
\[
 \max_x \sum_{i,C}\pi_i x_{iC}W_{iC} \tag{3}
\]
subject to $x_{iC}\geq0$, $\sum_Cx_{iC}=1$, and
\begin{align}
 \sum_Cx_{iC}U_{iC}&\geq0 &&\text{for all }i,\tag{4}\\
 \sum_Cx_{iC}U_{iC}&\geq\sum_Cx_{jC}U_{iC}
                         &&\text{for all }i,j,\tag{5}\\
 \sum_Cx_{iC}L_{iC}&\geq0 &&\text{for all }i,\tag{6}\\
 \sum_{i,C}\pi_i x_{iC}S_{iC}&\leq B. &&\tag{7}
\end{align}
For $B\geq0$ this linear program is feasible and attains a finite optimum.
It includes all effort deviations after every report, despite having only
linear reporting constraints.
\end{theorem}
\begin{proof}
The truthful payoff of type $i$ is the left side of (5). Following report
$j$, the contract is drawn from row $j$ and effort is then optimized,
so the deviation payoff is its right side. No additional deterministic
report can improve utility if (5) holds; a randomized report is a convex
combination of these payoffs and cannot improve either. Every effort
choice has already been maximized in the definition of $U_{iC}$.
Equations (4), (6), and (7) are exactly interim entrepreneur participation,
per-reported-type lender participation, and the common expected public
budget. Hence every allowed truthful mechanism gives a feasible matrix.
Conversely a feasible matrix can be implemented by the declared draw,
financing and effort timing; (4)--(7) verify each required constraint.
Its welfare is (3) by finite additivity of expectation.

All-rejection is feasible. Matrices lie in a finite product of simplices,
which is bounded and closed, and the additional inequalities are closed.
A maximizing sequence has a convergent subsequence in this set; the limit
is feasible and has the limiting linear objective. All coefficients are
finite, so the optimum is finite and attained.
\end{proof}

This reduction can compare a library of guarantee contracts with a library
including equity co-investment under the same budget. Enlarging the library
weakly increases the value, since an old feasible matrix extends by zero
probabilities. It does not imply that either instrument strictly dominates:
that conclusion depends on the library coefficients and incentive constraints.
The linear program is an exact finite-family characterization, not a claim
that a finite library approximates unrestricted contracts without error.

\section{Identification and unresolved parts}
Funding outcomes at a single contract reveal a selected success frequency,
not separately latent quality and effort. For instance an observed rate
$1/2$ can arise from effort zero with $p_0=1/2<p_1$, or effort one with
$p_0< p_1=1/2$, after choosing costs and success rents to satisfy the
corresponding strict effort inequality. The two structures imply different
responses to a changed repayment. Thus the coefficients of (3)--(7) need
causal contract variation or additional structural restrictions; the
program itself supplies no identification theorem.

The results leave continuous quality, endogenous screening, equilibrium
investor entry, spillover measurement, and general-equilibrium business
stealing unresolved. The weak-effort selector and interim commitment
convention are material restrictions. No empirical estimates or numerical
experiments are reported.

\begin{thebibliography}{9}
\bibitem{akcigit} U. Akcigit. \emph{Creative destruction in economic growth}.
The Scandinavian Journal of Economics 128(3), 509--541, 2026.
Section 7, item 1; publisher full text consulted 9 October 2026.
\url{https://onlinelibrary.wiley.com/doi/10.1111/sjoe.70037}.
\end{thebibliography}

\end{document}
